\documentclass[10pt,fleqn]{beamer}

\usepackage[portuges,brazilian]{babel}
\usepackage[T1]{fontenc}
\usepackage[utf8]{inputenc}

\usepackage{multirow} % To merge cels horizontaly

\usepackage{multicol} % para utilizar ambiente "multicols"
\usepackage{pdfpages} % para incluir paginas de documentos PDFs
\usepackage{color, amsmath}
\usepackage{subfigure}
\usepackage{graphicx} % para incluir figuras
\usepackage[ruled,longend]{algorithm2e} % para escrever algoritmos
\usepackage{ulem}       % para trocar a fonte?

\graphicspath{{pics/}}

\definecolor{defblue}{rgb}{0.2, 0.2, 0.7}
\definecolor{defred}{rgb}{0.4, 0.0, 0.0}
\definecolor{defgreen}{rgb}{0.0, 0.4, 0.0}

\definecolor{ninfagreen}{rgb}{0.168, 0.541, 0.569} % changed this

% Set left ident of equations
\setlength{\mathindent}{4pt}

\newcommand{\mydef}[1]{{\vspace{8pt} \large \color{defblue} #1}}
\newcommand{\myred}[1]{{\bf \color{defred} #1}}
\newcommand{\mygreen}[1]{{\bf \color{defgreen} #1}}

\renewcommand{\rmdefault}{put} % Utopia Font
\renewcommand{\sfdefault}{put} % Utopia Font
\renewcommand{\emph}[1]{ {\bf #1} }

% Scale Equations
\newcommand*{\Scale}[2][4]{\scalebox{#1}{$#2$}}%

\newcommand{\litem}[1]{
  \item{#1 \vspace{9pt}}
}

%\usetheme{Berkeley}
%\usetheme{Hannover}
%\usetheme{Marcos}
%\usetheme{default}
\usetheme{Boadilla}
%\usetheme{Pittsburgh}

%\useinnertheme{rectangles}
%\useoutertheme{tree}
%\setbeamercolor{separation line}{use=structure,bg=structure.fg!50!bg}

%\usecolortheme{dolphin}
%\usecolortheme{seahorse}

\beamertemplatenavigationsymbolsempty

% Escrita dos algoritmos
\SetKwInput{KwIn}{Entrada}
\SetKwInput{KwOut}{Saída}
\SetKwFor{While}{enquanto}{faça:}{fim}
\SetKwFor{For}{para}{faça:}{fim}
\SetKwBlock{Begin}{início}{fim}
\SetKwIF{If}{ElseIf}{Else}{se}{então}{senão se}{senão}{fim}

\newcommand{\bigbullet}{\,\begin{picture}(-1,1)(-1,-2)\circle*{5}\end{picture}\ }
\newcommand{\termo}[1]{\textbf{#1}}
\newcommand{\titulo}[1]{\centering \LARGE \vfill \textcolor{defblue}{\textbf{#1}} \vfill }

\newcommand{\chamada}[1]{
  \textcolor{ninfagreen}{\textbf{#1}}
  \vspace{10pt}
}

%\title{Um Estudo da Eficiência da Autocentralidade no Problema de Isomorfismo de Grafos}
\title[]{Uma Análise Experimental de Métodos para Classificação Multirrótulo}
\subtitle{}
\author{Lucas Henrique Sousa Mello}
\institute[NINFA]{}
\date[Vitória, 17/03/2014]{Vitória, 17 de Março de 2014}
\subject{}

% Figura de fundo
\setbeamercolor{background canvas}{bg=white}
\setbeamertemplate{background}{\includegraphics[width=\paperwidth,bb=0 0 1036 777]{fundo.png}}

% Posição dos títulos dos slides
\addtobeamertemplate{frametitle}{\vskip 2.5ex }{}
\addtobeamertemplate{frametitle}{\bf }{}

\begin{document}

\definecolor{beamer@blendedblue}{rgb}{0.168, 0.541, 0.569} % changed this

%\setbeamercolor{normal text}{fg=black,bg=white}
%\setbeamercolor{alerted text}{fg=red}
%\setbeamercolor{example text}{fg=green!50!black}

%\setbeamercolor{structure}{fg=beamer@blendedblue}

%\setbeamercolor{background canvas}{parent=normal text}
%\setbeamercolor{background}{parent=background canvas}

%\setbeamercolor{palette primary}{fg=yellow,bg=yellow} % changed this
%\setbeamercolor{palette secondary}{use=structure,fg=structure.fg!100!green} % changed this
%\setbeamercolor{palette tertiary}{use=structure,fg=structure.fg!100!green} % changed this

\begin{frame}
  \titlepage
\end{frame}

% \begin{frame}
%   \frametitle{Sumário}
%   \begin{itemize}
%     \litem{Introdução}
%     \litem{O problema da Mochila}
%     \litem{O Problema de Combate a Perdas}
%     \litem{Objetivos da Pesquisa}
%   \end{itemize}
% \end{frame}

\section{Introdução}

\begin{frame}
  \frametitle{Problemas de Classificação}
  \framesubtitle{O que são?}
  \begin{center}
  Determinar a classe de um objeto \\ dentre um conjunto pré-definido de classes. 
  \end{center}

  
  
  \begin{exampleblock}{Exemplos:}
  \begin{itemize}
   \item Fruta: descobrir se é Maçã, Uva ou Morango.
   \item Filme: descobrir se é de comédia ou não.
  \end{itemize}
  
 \end{exampleblock}
\end{frame}

\begin{frame}
  \frametitle{Problemas de Classificação}
  \framesubtitle{O que são?}
  
  Um objeto é descrito por um vetor de características.
  
    \begin{exampleblock}{Exemplo:}
  \begin{itemize}
   \item Fruta: Cor, massa e formato.
  \end{itemize}
  
 \end{exampleblock}
\end{frame}

\begin{frame}
  \frametitle{Problemas de Classificação Multirrótulo}
  \framesubtitle{O que são?}
  
  Um objeto pode estar associado a mais de uma classe simultaneamente.
  
    \begin{exampleblock}{Exemplo:}
  \begin{itemize}
   \item Filme: pode ser de comédia e romance.
  \end{itemize}
  
 \end{exampleblock}
\end{frame}


\begin{frame}
  \frametitle{Problemas de Classificação}
  \framesubtitle{O que são?}
  
  Definição geral de um problema de classificação unirrótulo:
  
  \vspace{10pt}
  \begin{align}
    \textrm{encontrar}
      \quad & f(x)=y \\
  	  \quad & \textrm{para todo }x \in X \nonumber \\
	  \nonumber
  \end{align}
  
  Muito difícil encontrar $f$, portanto $f$ é aproximado.
  
\end{frame}


\begin{frame}
  \frametitle{Problemas de Classificação}
  \framesubtitle{O que são?}
  
  Definição geral de um problema de classificação multi:
  
  \vspace{10pt}
  \begin{align}
    \textrm{encontrar}
      \quad & f(x)=y \\
  	  \quad & \textrm{para todo }x \in X \nonumber \\
	  \nonumber
  \end{align}
  
  Muito difícil encontrar $f$, portanto $f$ é aproximado.
  
\end{frame}

%\begin{frame}
%  \frametitle{Otimização Contínua x Combinatória (Discreta)}
%  \begin{align*}
%    \textrm{maximizar}
%      \quad & z = f(x) & \textrm{maximizar} \quad & z = f(x) \\
%  	  \quad & x \in \mathbb{R} & \quad & x \in \mathbb{N} \\
%	  \nonumber
%  \end{align*}
%\end{frame}

\begin{frame}
  \frametitle{Problemas de Otimização Combinatória}
  \framesubtitle{Quando o problema é combinatório?}
  Definição geral de um problema de otimização combinatória:
  \begin{align}
    \textrm{maximizar}
      \quad & z = f(x) \nonumber \\
  	  \quad & x \in S \nonumber, \qquad S \, \textrm{é finito} \\
	  \nonumber
  \end{align}
\end{frame}

%\begin{frame}
%  \frametitle{Otimização Contínua x Combinatória (Discreta)}
%  \begin{align*}
%    \textrm{maximizar}
%      \quad & z = f(x) & \textrm{maximizar} \quad & z = f(x) \\
%  	  \quad & x \in \mathbb{R} & \quad & x \in \mathbb{N} \\
%	  \nonumber
%  \end{align*}
%\end{frame}

%\begin{frame}
%  \frametitle{Problemas Difíceis (NP-Hard)}
%\end{frame}

\begin{frame}
  \frametitle{Problemas de Otimização Combinatória}
  Muitos deles são NP-Hard. Como resolvê-los?
  \vspace{10pt}
  \begin{itemize}
    \litem{Algoritmos Exatos}
    \litem{Algoritmos de Aproximação}
    \litem{Algoritmos Heurísticos}
	\begin{itemize}
	  \item{Meta-heurísticas}
	\end{itemize}
  \end{itemize}
\end{frame}

% \section{O problema da Mochila}
% 
% \begin{frame}
%   \frametitle{O Problema da Mochila} % versão padrão
%   \framesubtitle{\textit{Knapsack Problem}}
%   Utilizar a mochila para carregar o maior valor possível, respeitando a sua
%   capacidade.
%   \begin{columns}[T]
%     % Primeira coluna
% 	  \begin{column}{.5\textwidth}
%       \begin{block}{Modelagem matemática}
% 	    { \footnotesize
%         \begin{align*}
%           \textrm{maximizar} \quad
% 		  	& z = \sum_{j=1}^n p_j x_j \\
%           \textrm{sujeito a} \quad
% 		  	& \sum_{j=1}^n w_{j}x_j \leq C \\
%             & {x_j \in \{0,1\}, \quad j=1,\cdots,n.}
%         \end{align*}
% 		}
%       \end{block}
%     \end{column}
%     % Segunda coluna
%     \begin{column}{.5\textwidth}
% 	  \begin{center}
%         \includegraphics[scale=0.3]{knapsack}
% 	  \end{center}
%     \end{column}
%   \end{columns}
% \end{frame}
% 
% \begin{frame}
%   \frametitle{O Problema da Mochila Multidimensional}
%   A \textit{mochila} agora possui mais de uma capacidade (dimensão).
%   \begin{columns}[T]
%     % Primeira coluna
% 	  \begin{column}{.5\textwidth}
%       \begin{block}{Modelagem matemática}
% 	    { \footnotesize
%         \begin{align*}
%           \textrm{maximizar}
%             \quad & z = \sum_{j=1}^n p_j x_j \\
%           \textrm{sujeito a}
%   	        \quad & \sum_{j=1}^n w_{ij}x_j \leq c_i, \quad i=1,\cdots,m, \\
%                 & x_j \in \{0,1\}, \quad j=1,\cdots,n.
%         \end{align*}
% 		}
%       \end{block}
%     \end{column}
%       % Segunda coluna
%       \begin{column}{.5\textwidth}
% 	  { \tiny
% 	  \begin{center}
% 	  \chamada{\small Ex.: Seleção de projetos.}
% 	  \\
%       \begin{tabular}{|l|r|c|c|c|} \hline
%       {\bf Projeto} & {\bf Valor} & {\bf Analistas} & {\bf Prog.} & {\bf Designers} \\ \hline
% 	  Projeto A & R\$ 100 000,00 & 2 & 6 & 1 \\ \hline
% 	  Projeto B & R\$ 150 000,00 & 4 & 10 & 3 \\ \hline
% 	  Projeto C & R\$ 80 000,00 & 2 & 4 & 2 \\ \hline
%       \end{tabular}
% 	  \end{center}
% 
% 	  \begin{center}
%       \begin{tabular}{|l|c|c|c|c|} \hline
%         {\bf Equipe} & {\bf Quantidade} \\ \hline
% 	    Analistas & 6 \\ \hline
% 	    Programadores & 16 \\ \hline
% 	    Designers & 4 \\ \hline
%       \end{tabular}
% 	  \end{center}
% 	  }
%     \end{column}
%   \end{columns}
% \end{frame}
% 
% \begin{frame}
%   \frametitle{O Problema da Mochila com Dependência}
%   Nesta variação existe uma dependência entre os itens selecionados. \\
%   \vfill
%   \begin{tabular}{|l|r|c|c|c|} \hline
%     {\bf Projeto} & {\bf Valor} & {\bf Analistas} & {\bf Programadores} & {\bf Designers} \\ \hline
% 	Projeto A & R\$ 100 000,00 & 2 & 6 & 1 \\ \hline
% 	Projeto B & R\$ 150 000,00 & 4 & 10 & 3 \\ \hline
% 	Projeto C & R\$ 80 000,00 & 2 & 4 & 2 \\ \hline
%   \end{tabular}
%   \vfill
%   \begin{columns}[T]
%     \begin{column}{.5\textwidth}
%       \begin{tabular}{|l|c|c|c|c|} \hline
%         {\bf Equipe} & {\bf Quantidade} \\ \hline
% 	    Analistas & 6 \\ \hline
% 	    Programadores & 16 \\ \hline
% 	    Designers & 4 \\ \hline
%       \end{tabular}
% 	\end{column}
%     \begin{column}{.5\textwidth}
% 	  \chamada{Ex.:} Alocando no mínimo $1$ Analista para cada $3$ programadores.
% 	\end{column}
%   \end{columns}
% \end{frame}
% 
% \section{O Problema de Combate a Perdas}
% 
% \begin{frame}
%   \frametitle{O Problema de Planejamento de Combate a Perdas}
%   \chamada{Contexto do problema}
%   \begin{itemize}
%   	\litem{A distribuidora local estima que do total de energia distribuida
% 	  $13\%$ é perdida com fraude e apenas $7\%$ com fenômenos físicos.}
%   	\litem{A ANEEL estipula para as ditribuidoras de todo o país uma meta anual de
%   	  redução de perda de energia.}
%     \litem {Planejamento plurianual para redução da fraude:}
% 	\begin{enumerate}
% 	  \item{Alcançar a meta estipulada, resgatando o maior valor em dinheiro possível;}
% 	  \item{Conjunto de ações de redução de fraude;}
% 	  \item{Respeitando um orçamento anual;}
% 	\end{enumerate}
%   \end{itemize}
% \end{frame}
% 
% \begin{frame}
%   \frametitle{O Problema de Planejamento de Combate a Perdas}
%     \begin{center}
% 	\includegraphics[scale=0.37]{meta.eps}
% 	\vspace{10pt}
%     { \scriptsize
%     \begin{tabular}{|l|r|r|r|r|r|r|r|}
%       \hline
%       \multirow{2}{*}{\bf Ação} & \multirow{2}{*}{\bf Quant.$^{1}$} & \multicolumn{2}{|c|}{\bf Custo} & \multicolumn{4}{c|}{\bf Recuperação (Kwh)} \\ \cline {3-4} \cline{5-8}
% 		& & {\bf CAPEX} & {\bf OPEX} & ${ \bf 1^\circ}$ & ${ \bf 2^\circ}$ & ${ \bf 3^\circ}$ & ${ \bf 4^\circ}$ \\ \hline
% 		Ação A & 140 & R\$ 100,00 & R\$ 0,00 & 100 & 70 & 70 & 70 \\ \hline
% 		Ação B & 200 & R\$ 50,00 & R\$ 70,00 & 20 & 50 & 75 & 50 \\ \hline
% 		Ação C & 175 & R\$ 90,00 & R\$ 30,00 & 50 & 100 & 100 & 100 \\ \hline
% 		Ação D & 1000 & R\$ 1050,00 &R\$ 250,00 & 13 & 600 & 600 & 600 \\ \hline
%     \end{tabular}
% 	}
%     \end{center}
% 	{\footnotesize $^1$Restrição de mercado.}
% \end{frame}
% 
% \begin{frame}
%   \frametitle{O Problema de Planejamento de Combate a Perdas}
%   \chamada{Outras características do problema:}
%   \begin{enumerate}
%     \item{Diferentes faixas de tensão (baixa, média ou alta);}
%     \begin{itemize}
%       \litem{Preço da energia varia entre as faixas.}
%     \end{itemize}
%     \litem{Deve-se considerar a taxa interna de retorno (juros);}
%     \litem{Pode haver dependência entre as ações;}
%     \litem{Mesmo com orçamento anual, é preferível realizar um planejamento por
% 	  períodos menores (bimestral, trimestral, mensal, etc);}
%   \end{enumerate}
% \end{frame}
% 
% \section{Objetivo da Pesquisa}
% 
% \begin{frame}
%   \frametitle{Objetivos da Linha de Pesquisa}
%   \begin{itemize}
%     \item{Estudar complexidade do problema;}
% 	\begin{itemize}
% 	  \item{Quais conjuntos de ações tornam o problema difícil?}
% 	  \litem{O quanto cada restrição torna o problema difícil?}
% 	\end{itemize}
%     \litem{Propor métodos heurísticos para a resolução.}
%   \end{itemize}
% \end{frame}


\end{document}
